% Pollis New Methods: panel econometrics (random-effects variance estimators) — draft 1 (11 October 2026)
% Compile with: latexmk -pdf nm_panel_econometrics.tex && latexmk -c
% Shared look for the Pollis New Methods documents (same fonts and colours as the brochure and the workflow handouts).
\documentclass[10pt,a4paper]{article}

\usepackage[a4paper,left=18mm,right=18mm,top=15mm,bottom=17mm,footskip=8mm]{geometry}
\usepackage[T1]{fontenc}
\usepackage[default,semibold]{sourcesanspro}
\usepackage{textcomp}
\usepackage[italic]{mathastext}
\usepackage{amsmath}
\usepackage{microtype}
\usepackage[table]{xcolor}
\usepackage{tikz}
\usepackage[most]{tcolorbox}
\usepackage{fontawesome5}
\usepackage{tabularx,booktabs}
\usepackage{enumitem}
\usepackage{fancyhdr}
\usepackage[hidelinks]{hyperref}

% ---------------------------------------------------------------- colours (as in the brochure)
\definecolor{navy}{HTML}{1C3A5E}
\definecolor{pblue}{HTML}{3B78B5}
\definecolor{amber}{HTML}{F2B33D}
\definecolor{deepamber}{HTML}{D9902A}
\definecolor{pale}{HTML}{EEF3F9}
\definecolor{ink}{HTML}{2A2F36}
\definecolor{logonavy}{HTML}{173D63}
\definecolor{logodark}{HTML}{C98A1E}
\definecolor{logoamber}{HTML}{F2B446}
\definecolor{logocream}{HTML}{FEEFD2}
\definecolor{soft}{HTML}{5E6773}
\definecolor{agree}{HTML}{2F7D4F}
\definecolor{differ}{HTML}{B23B3B}
\color{ink}

\setlength{\parindent}{0pt}
\setlength{\parskip}{4pt}
\setlist{nosep,leftmargin=*,itemsep=2pt}
\renewcommand{\arraystretch}{1.15}

\pagestyle{fancy}
\fancyhf{}
\renewcommand{\headrulewidth}{0pt}
\fancyfoot[L]{\footnotesize\color{soft}\textbf{Pollis}\enspace New Methods\enspace·\enspace\nmtheme}
\fancyfoot[R]{\footnotesize\color{soft}\thepage}
\newcommand{\nmtheme}{}
\newcommand{\nmdate}{11 October 2026}

% The Pollis logo, redrawn so it prints sharply at any size
\newcommand{\pollislogo}[1]{%
  \begin{tikzpicture}[x=#1/512,y=-#1/512,baseline=(current bounding box.center)]
    \fill[logonavy] (0,0) rectangle (512,512);
    \fill[logodark] (36,422) -- (104,256) -- (172,422) -- cycle;
    \fill[logoamber] (104,422) -- (203,180) -- (302,422) -- cycle;
    \fill[logocream] (203,422) -- (339,89) -- (476,422) -- cycle;
  \end{tikzpicture}}

% Top of the first page: logo, name; then the kicker, title and a one-line promise
\newcommand{\nmhead}[3]{%
  \renewcommand{\nmtheme}{#1}%
  \thispagestyle{fancy}%
  \begin{minipage}[c]{13mm}\pollislogo{11mm}\end{minipage}%
  \begin{minipage}[c]{\dimexpr\linewidth-13mm\relax}
    {\color{navy}\Large\bfseries Pollis}\enspace{\color{soft}New Methods}\hfill
    {\color{deepamber}\small\bfseries\MakeUppercase{#1}}
  \end{minipage}\par\vspace{7mm}
  {\color{navy}\fontsize{22}{26}\selectfont\bfseries #2}\par\vspace{4pt}
  {\color{soft}\large #3}\par\vspace{5mm}}

\newcommand{\nmsection}[1]{\par\vspace{5pt}{\color{navy}\large\bfseries #1}\par\vspace{1pt}}
\newcommand{\nmsub}[1]{\par\vspace{2pt}{\color{deepamber}\bfseries #1}\par}
\newcommand{\gt}{\,{\color{deepamber}\bfseries ›}\,}
\newcommand{\menupath}[1]{{\bfseries\color{navy}#1}}
\newcommand{\fn}[1]{\texttt{\small #1}}
% Citations: \nmcite{key}{Breiman (2001)} links to the entry \nmref{key}{...} in the References
\newcommand{\nmcite}[2]{\hyperlink{ref:#1}{\color{pblue}#2}}
\newcommand{\nmref}[2]{\raisebox{1.1em}[0pt][0pt]{\hypertarget{ref:#1}{}}#2\par}
\newcommand{\nmdoi}[1]{\href{https://doi.org/#1}{\color{pblue}doi:#1}}
% Where a notebook is in Pollis: its title in Browse Methods, and its file
\newcommand{\browse}[2]{\menupath{Compose \gt New Methods \gt Browse Methods}: \emph{#1} (\fn{#2.ipynb})}

% Key facts at the top: source, missing in Julia, notebook, result
\newcommand{\factsheet}[4]{%
  \begin{tcolorbox}[enhanced, colback=pale, colframe=pale, boxrule=0pt, sharp corners,
    left=7pt, right=7pt, top=5pt, bottom=5pt, fontupper=\small]
  \begin{tabularx}{\linewidth}{@{}>{\leavevmode\bfseries\color{navy}}p{27mm}>{\leavevmode}X@{}}
  Tested against & #1\\
  Missing in Julia & #2\\
  Notebook & #3\\
  Result & #4
  \end{tabularx}
  \end{tcolorbox}}

% Boxes
\newtcolorbox{pbox}[2][]{enhanced, colback=pale, colframe=amber, boxrule=0pt,
  leftrule=3pt, sharp corners, left=7pt, right=7pt, top=6pt, bottom=6pt,
  fontupper=\small\raggedright, before upper={{\color{navy}\normalsize\bfseries #2}\par\smallskip}, #1}
\newtcolorbox{nbox}[2][]{enhanced, colback=navy, colframe=navy, boxrule=0pt,
  sharp corners, left=8pt, right=8pt, top=7pt, bottom=7pt, coltext=white,
  fontupper=\small\raggedright, before upper={{\color{amber}\normalsize\bfseries #2}\par\smallskip}, #1}

% Value tables: published beside Julia
\newenvironment{valuetable}[1]{%
  \par\small\tabularx{\linewidth}{#1}\toprule}{\bottomrule\endtabularx\par}
\newcommand{\ok}{{\color{agree}\faCheck}}
\newcommand{\no}{{\color{differ}\faTimes}}

% End of the document: credit and licence
\newcommand{\nmcredit}{%
  \par\vfill
  \begin{minipage}[b]{0.78\linewidth}\footnotesize\color{soft}
  Methods written and tested with Claude Code (Anthropic) in Pollis, from the published papers and their data.\par
  \textcopyright\ 2026 Antonio Saragga Seabra\enspace·\enspace Licensed under \href{https://creativecommons.org/licenses/by-sa/4.0/}{CC BY-SA 4.0}\enspace·\enspace Draft, \nmdate
  \end{minipage}\hfill
  \begin{minipage}[b]{0.18\linewidth}\raggedleft\pollislogo{9mm}\par\vspace{2pt}{\footnotesize contact@pollis.org}\end{minipage}}


\begin{document}

\nmhead{Panel Econometrics}
  {Random effects: WALHUS, AMEMIYA, Swamy-Arora and the nested model}
  {Five ways to split a panel's error variance, written from one rule and tested against the tables of \nmcite{baltagi2005}{Baltagi (2005)}}

\factsheet
  {\nmcite{baltagi2005}{Baltagi (2005)}, \emph{Econometric Analysis of Panel Data}, 3rd ed.: Tables 2.1, 2.5, 2.6, 3.1, 3.6, 3.7 and 9.4.
   The investment data of \nmcite{grunfeld1958}{Grunfeld (1958)}, the OECD gasoline data of \nmcite{baltagigriffin1983}{Baltagi and Griffin (1983)} and the public capital data of \nmcite{munnell1990}{Munnell (1990)} (PollisDatasets).}
  {Wallace-Hussain (WALHUS), Amemiya (AMEMIYA), the two-way Swamy-Arora, and the estimators of the nested model
   (states within regions). Julia packages offer the one-way Swamy-Arora at most; R's plm has the rest.}
  {\browse{WALHUS, AMEMIYA, Swamy-Arora and nested error components}{panelre}}
  {288 of 298 printed values to every digit. R's plm agrees with Julia to its 7 printed significant digits wherever
   both use the same convention.}

\nmsection{The problem}
A random-effects panel regression is a GLS regression whose weights depend on the variances of the error components:
an individual effect, perhaps a time effect or a nested effect, and the remainder. The estimators differ only in how they
estimate those variances before GLS, and the choice moves the standard errors and sometimes the coefficients. A textbook
that prints a row for each estimator cannot be followed in a language that has one of them.

\nmsection{One rule for all the estimators}
Let $e = My$ be residuals, where $M$ removes the regressors, so $e = Mu$ whatever the coefficients. A sum of squares
$q = e'Ae$ has the exact expectation
\[ \mathrm{E}(q) = \textstyle\sum_c \sigma_c^2 \,\mathrm{tr}(M'AM\,Z_cZ_c'), \]
with $Z_c$ the incidence matrix of component $c$. With as many sums of squares as components, equating each to its
expectation gives the variances. This is how \nmcite{baltagi2005}{Baltagi (2005)} defines the unbalanced estimators (Section 9.2.1, eqs.\ 9.16--9.21),
and in balanced panels it gives the familiar formulas. The estimators are then only choices of residuals:
\begin{itemize}
  \item \textbf{WALHUS}, \nmcite{wallace1969}{Wallace and Hussain (1969)}: OLS residuals in every sum of squares.
  \item \textbf{AMEMIYA}, \nmcite{amemiya1971}{Amemiya (1971)}, and \nmcite{wansbeek1989}{Wansbeek and Kapteyn (1989)} when unbalanced: the residuals of the Within regression.
  \item \textbf{SWAR}, \nmcite{swamy1972}{Swamy and Arora (1972)}: each sum of squares from its own regression, Within or Between.
  \item \textbf{HFB}, Henderson's method III, \nmcite{henderson1953}{Henderson (1953)}, as in \nmcite{fuller1974}{Fuller and Battese (1974)}: reductions in the sum of squares as the nested dummies are added.
  \item \textbf{MIVQUE}, \nmcite{rao1971}{Rao (1971)}, with a prior guess of the variances.
\end{itemize}
A negative variance is set to zero, as in Tables 3.1 and 3.7. Writing the rule once, rather than five sets of formulas,
is what made the nested, unbalanced case as easy as the one-way case.

\nmsection{What the tests show}
\begin{valuetable}{@{}lXl@{}}
\textbf{Table} & \textbf{Model and data} & \textbf{Agreement}\\\midrule
2.1, p.\,21 & Grunfeld, one-way: all estimators & every printed digit\\
2.5, p.\,24 & Gasoline, one-way: WALHUS, AMEMIYA & every printed digit\\
2.5, p.\,24 & Gasoline, one-way: SWAR & variances agree; price $-.420$ vs $-.402$ (see below)\\
2.6, p.\,25 & Public capital, one-way & every printed digit\\
3.1, p.\,43 & Grunfeld, two-way & every printed digit\\
3.6, p.\,46 & Gasoline, two-way & every printed digit\\
3.7, p.\,47 & Public capital, two-way & every printed digit\\
9.4, p.\,182 & Public capital, nested: WH, WK, SA, MV1, MV2 & every printed digit\\
9.4, p.\,182 & nested: HFB & six values in the last digit\\
\end{valuetable}
The MLE and REML columns of Table 9.4 also agree to every printed digit; they are fitted with MixedModels.jl and check
the data and the model. The EViews output Baltagi reprints for WALHUS in Table 2.2 is matched to every digit EViews prints:
$\sigma_\mu = 87.358027$ and $\sigma_\nu = 53.745175$ here, against 87.35803 and 53.74518.

\nmsection{What the replication found}
\nmsub{The book uses two conventions, and both had to be matched}
The standard errors of the two-way tables (3.1, 3.6, 3.7) are scaled by the residual variance of the transformed
regression, as EViews does; those of the one-way tables and Table 9.4 are not. With one convention for all, about a third
of the standard errors differ in the last digit. And the sums of squares are equated to their exact expectations, not
divided by the traces of eqs.\ (2.21)--(2.22): this gives AMEMIYA's remainder SWAR's divisor $N(T-1)-K$, which is why the two
share $\sigma_\nu = 52.77$ in Table 2.1.

\nmsub{R's plm does not reproduce the book's WALHUS and AMEMIYA}
plm divides by traces, so for one-way WALHUS it gives $\sigma_\mu = .078$ against Baltagi's .082 on the public capital data,
and 75.43 against 87.36 on Grunfeld. In the two-way models it divides the between sums of squares by $N$ and $T$ rather than
$N-1$ and $T-1$. Julia reproduces plm exactly with plm's conventions (\fn{divisors = :trace}, \fn{scale = true}), and the
book exactly with the exact expectations: the two programs differ by a definition, not by an error.

\nmsub{A probable misprint in Table 2.5}
For SWAR on the gasoline data, the variances agree ($\rho = .82$, $\sigma_\mu = .196$, $\sigma_\nu = .092$), but the price
coefficient is $-.420$ here against $-.402$ in print. R's plm gives $-.4204$, as Julia does. The printed value looks like
transposed digits.

\nmsection{Still open}
\begin{itemize}
  \item \textbf{One standard error in Table 2.5.} With the two-way scaling, Julia's SWAR standard errors are .059, .040 and
    .026 against .059, .042 and .026; plm gives .059, .040 and .0255. The .042 is not explained.
  \item \textbf{HFB in Table 9.4}, in the last digit: intercept 2.088 vs 2.084, log K .271 vs .272, log L .744 vs .743,
    its standard error .151 vs .150, $\sigma_\mu^2$ .0027 vs .0029, $\sigma_\nu^2$ .0046 vs .0044. The fitting-constants
    reductions follow Section 9.2.1; \nmcite{bsj2001}{Baltagi, Song and Jung (2001)} define the nested version, which may order them differently.
  \item \textbf{MV3.} The book does not give the MIVQUE priors. Priors $(0,0,1)$ and $(1,1,1)$ reproduce MV1 and MV2
    exactly; no prior tried reproduces MV3 (the closest gives $\sigma_\mu^2$ .0018 against .0017), so MV3 is not tested.
\end{itemize}

\nmsection{Further considerations}
\begin{itemize}
  \item \textbf{A note for users of panel software.} Two programs can both be ``right'' and still disagree in the second
    significant digit: the divisor and the scaling of the standard errors are conventions, and a replication has to name them.
    The notebook makes both explicit options, so either convention can be chosen deliberately.
  \item \textbf{Upstream.} The single quadratic-form rule is small and general enough to be offered to a Julia panel-data
    package, with these tables as its tests.
  \item \textbf{Next.} The \nmcite{bsj2001}{Baltagi, Song and Jung (2001)} nested HFB, to close the last-digit gap; the unbalanced two-way model;
    and the estimator of \nmcite{hausman1981}{Hausman and Taylor (1981)}, also in Baltagi's chapters and also missing in Julia.
\end{itemize}

\nmsection{References}
{\footnotesize
\nmref{amemiya1971}{Amemiya, T. (1971). The estimation of the variances in a variance-components model. \emph{International Economic Review} 12(1), 1--13.}
\nmref{baltagi2005}{Baltagi, B.\,H. (2005). \emph{Econometric Analysis of Panel Data}, 3rd ed. Wiley.}
\nmref{baltagigriffin1983}{Baltagi, B.\,H. and Griffin, J.\,M. (1983). Gasoline demand in the OECD: an application of pooling and testing procedures. \emph{European Economic Review} 22(2), 117--137.}
\nmref{bsj2001}{Baltagi, B.\,H., Song, S.\,H. and Jung, B.\,C. (2001). The unbalanced nested error component regression model. \emph{Journal of Econometrics} 101(2), 357--381. \nmdoi{10.1016/S0304-4076(00)00089-0}}
\nmref{fuller1974}{Fuller, W.\,A. and Battese, G.\,E. (1974). Estimation of linear models with crossed-error structure. \emph{Journal of Econometrics} 2(1), 67--78.}
\nmref{grunfeld1958}{Grunfeld, Y. (1958). \emph{The Determinants of Corporate Investment}. PhD thesis, University of Chicago.}
\nmref{hausman1981}{Hausman, J.\,A. and Taylor, W.\,E. (1981). Panel data and unobservable individual effects. \emph{Econometrica} 49(6), 1377--1398.}
\nmref{henderson1953}{Henderson, C.\,R. (1953). Estimation of variance and covariance components. \emph{Biometrics} 9(2), 226--252.}
\nmref{munnell1990}{Munnell, A.\,H. (1990). How does public infrastructure affect regional economic performance? \emph{New England Economic Review}, September/October, 11--32.}
\nmref{rao1971}{Rao, C.\,R. (1971). Minimum variance quadratic unbiased estimation of variance components. \emph{Journal of Multivariate Analysis} 1(4), 445--456.}
\nmref{swamy1972}{Swamy, P.\,A.\,V.\,B. and Arora, S.\,S. (1972). The exact finite sample properties of the estimators of coefficients in the error components regression models. \emph{Econometrica} 40(2), 261--275.}
\nmref{wallace1969}{Wallace, T.\,D. and Hussain, A. (1969). The use of error components models in combining cross section with time series data. \emph{Econometrica} 37(1), 55--72.}
\nmref{wansbeek1989}{Wansbeek, T. and Kapteyn, A. (1989). Estimation of the error-components model with incomplete panels. \emph{Journal of Econometrics} 41(3), 341--361.}
}

\nmcredit
\end{document}
